\documentclass[12pt]{ximera}
\title{Homework 5: Shapley-Shubik Power, Subsets, and Permutations}
\begin{document}
\begin{abstract}
Sections 2.3--2.4: pivotal players in sequential coalitions and the Shapley-Shubik
power distribution, counting coalitions and sequential coalitions, and using coalition
\emph{types} to handle a system too large to enumerate by hand.
\end{abstract}
\maketitle

\section*{Sections 2.3--2.4: Shapley-Shubik Power}

In a \emph{sequential} coalition, players join one at a time. The \emph{pivotal}
player is the one whose arrival first brings the running total up to the quota.

\begin{problem}
Consider the weighted voting system $[55:40,30,20,10]$ with players $A$, $B$, $C$, and
$D$. There are $4! = 24$ sequential coalitions.

\begin{prompt}
\textbf{(a)} How many of the $24$ sequential coalitions have each player as the
pivotal player?

$A$: $\answer{10}$ \qquad $B$: $\answer{6}$ \qquad $C$: $\answer{6}$ \qquad $D$: $\answer{2}$

\begin{hint}
Take one sequential coalition at a time and keep a running total, stopping at the
first player who brings it to $55$ or more. For example, in
$\langle D, B, A, C\rangle$ the running totals are $10$, $40$, $80$ --- so $A$ is
pivotal.
\end{hint}

\begin{feedback}
Working through all $24$ orderings and marking where the running total first reaches
$55$:
\begin{itemize}
\item $A$ is pivotal in $10$ of them,
\item $B$ in $6$,
\item $C$ in $6$,
\item $D$ in $2$.
\end{itemize}
These total $24$, exactly one pivotal player per sequential coalition. \checkmark

$D$ is pivotal only in the two orderings $\langle B,C,D,A\rangle$ and
$\langle C,B,D,A\rangle$ --- the cases where $B$ and $C$ arrive first (reaching $50$,
just short) and $D$'s $10$ votes push the total to $60$.
\end{feedback}
\end{prompt}

\begin{prompt}
\textbf{(b)} Determine the Shapley-Shubik power distribution for this weighted voting
system, as percentages rounded to one decimal place.

$A$: $\answer[tolerance=0.05]{41.7}\%$ \qquad
$B$: $\answer[tolerance=0.05]{25}\%$ \qquad
$C$: $\answer[tolerance=0.05]{25}\%$ \qquad
$D$: $\answer[tolerance=0.05]{8.3}\%$

\begin{feedback}
Divide each pivotal count by the total of $24$:
\[
A: \frac{10}{24}\approx 41.7\%, \quad
B: \frac{6}{24}=25\%, \quad
C: \frac{6}{24}=25\%, \quad
D: \frac{2}{24}\approx 8.3\%.
\]

Note that $B$ and $C$ have equal power ($25\%$ each) despite $B$ having $30$ votes and
$C$ only $20$. And $D$, with $10$ votes --- a quarter of $A$'s $40$ --- holds only
$8.3\%$ of the power against $A$'s $41.7\%$, a fifth as much. Power is not
proportional to weight.
\end{feedback}
\end{prompt}
\end{problem}

\begin{problem}
Consider a weighted voting system with $12$ players ($P_1,\ldots,P_{12}$).

\begin{prompt}
\textbf{(a)} What is the \emph{total} number of possible coalitions (subsets of
players) in this system?

$\answer{4096}$

\begin{feedback}
Each of the $12$ players is independently either in or out, so there are
\[ 2^{12} = 4096 \]
subsets, counting the empty one.
\end{feedback}
\end{prompt}

\begin{prompt}
\textbf{(b)} What is the number of possible \emph{nonempty} coalitions?

$\answer{4095}$

\begin{feedback}
Discard the empty coalition: $2^{12}-1 = 4095$.
\end{feedback}
\end{prompt}

\begin{prompt}
\textbf{(c)} How many coalitions include $P_1$ and $P_2$ but not $P_3$?

$\answer{512}$

\begin{hint}
Three of the twelve players have their membership already decided. How many players
are still free to be in or out?
\end{hint}

\begin{feedback}
The membership of $P_1$, $P_2$, and $P_3$ is fixed, leaving the other $9$ players free:
\[ 2^9 = 512. \]
\end{feedback}
\end{prompt}

\begin{prompt}
\textbf{(d)} What is the total number of \emph{sequential} coalitions in this system?

$\answer{479001600}$

\begin{feedback}
A sequential coalition is an ordering of all $12$ players, so there are
\[ 12! = 479{,}001{,}600 \]
of them. This is why coalition \emph{types} matter --- enumerating half a billion
orderings by hand is not an option.
\end{feedback}
\end{prompt}
\end{problem}

\begin{problem}
True story: when I was young, my household regularly voted on things like where to eat
using a weighted voting system where each child had one vote, each adult had two
votes, and a simple majority was sufficient to make a decision.

For a household with two adults and three children, this system has $5!$ sequential
coalitions, but we can use the shorthand \textbf{A} for `adult' and \textbf{C} for
`child', since players within those roles are indistinguishable.

\begin{prompt}
\textbf{(a)} How many sequential coalition \emph{types} are there --- that is, how many
ways are there to arrange two \textbf{A}'s and three \textbf{C}'s in a sequence?

$\answer{10}$

\begin{feedback}
Choose which $2$ of the $5$ positions hold the \textbf{A}'s:
\[ C(5,2) = 10 \]
arrangements.
\end{feedback}
\end{prompt}

\begin{prompt}
\textbf{(b)} How many distinct sequential coalitions are there of each type, and how
many are there altogether?

Each type represents $\answer{12}$ distinct sequential coalitions, for a total of
$\answer{120}$.

\begin{hint}
Within a type, the two adults could be in either order, and the three children in any
order.
\end{hint}

\begin{feedback}
Within a given type, the two adults can be arranged in $2!=2$ ways and the three
children in $3!=6$ ways, so each type covers
\[ 2!\cdot 3! = 12 \]
distinct sequential coalitions. Altogether $10 \times 12 = 120 = 5!$. \checkmark

The total weight is $2+2+1+1+1 = 7$, so a simple majority means a quota of $q = 4$.
Marking the pivotal role in each type:
\[
\begin{array}{ll}
\langle A,A,C,C,C\rangle \to \textbf{A} &
\langle C,A,C,A,C\rangle \to \textbf{C}\\
\langle A,C,A,C,C\rangle \to \textbf{A} &
\langle C,A,C,C,A\rangle \to \textbf{C}\\
\langle A,C,C,A,C\rangle \to \textbf{C} &
\langle C,C,A,A,C\rangle \to \textbf{A}\\
\langle A,C,C,C,A\rangle \to \textbf{C} &
\langle C,C,A,C,A\rangle \to \textbf{A}\\
\langle C,A,A,C,C\rangle \to \textbf{A} &
\langle C,C,C,A,A\rangle \to \textbf{A}
\end{array}
\]
So an adult is pivotal in $6$ of the $10$ types, and a child in $4$.
\end{feedback}
\end{prompt}

\begin{prompt}
\textbf{(c)} What is the pivotal count for an individual adult in this system?

$\answer{36}$

\begin{feedback}
Adults are pivotal in $6$ types, each covering $12$ distinct sequential coalitions,
so the two adults are pivotal $6\times 12 = 72$ times between them. Split evenly:
\[ 72 / 2 = 36 \]
per adult.
\end{feedback}
\end{prompt}

\begin{prompt}
\textbf{(d)} What is the pivotal count for an individual child?

$\answer{16}$

\begin{feedback}
Children are pivotal in $4$ types, so $4\times 12 = 48$ sequential coalitions have a
child pivotal. Split among the three children:
\[ 48/3 = 16 \]
each. As a check, $2(36) + 3(16) = 72 + 48 = 120$. \checkmark
\end{feedback}
\end{prompt}

\begin{prompt}
\textbf{(e)} Determine the Shapley-Shubik power distribution for this five-player
system, as percentages rounded to two decimal places.

Each adult: $\answer[tolerance=0.005]{30}\%$ \qquad
Each child: $\answer[tolerance=0.005]{13.33}\%$

\begin{feedback}
\[
\text{Adult}: \frac{36}{120}=30\%, \qquad \text{Child}: \frac{16}{120}\approx 13.33\%.
\]
Check: $2(30\%) + 3(13.33\%) = 60\% + 40\% = 100\%$. \checkmark

An adult has twice a child's votes but well over twice a child's power: $30\%$ versus
$13.33\%$, a ratio of $2.25$ to $1$. The two parents together hold $60\%$ of the
power with only $4$ of the $7$ votes.
\end{feedback}
\end{prompt}
\end{problem}

\end{document}
